\documentclass[10pt,a4paper]{amsart}
\usepackage[margin=19mm]{geometry}
\usepackage{amsmath,amssymb,mathtools}
\usepackage[scr=rsfs,cal=euler]{mathalfa}
\usepackage{xfrac}
\usepackage[unicode,hidelinks,bookmarks=false]{hyperref}
\newcommand{\ie}{i.e.\ }
\newcommand{\cf}{cf.\ }
\newcommand{\resp}{respectively\ }
\newcommand{\Subsection}{\subsection}
\providecommand{\nobreakdash}{\nobreak\hbox{-}}
\setlength{\emergencystretch}{2em}
\multlinegap0pt
\usepackage[shortlabels]{enumitem}
\usepackage{amssymb}
\usepackage{bm}
\usepackage{algorithm}\usepackage{algpseudocode}
\usepackage{xcolor}
\newtheorem{theorem}{Theorem}
\title{Deep-Koopman-KANDy: Dictionary Discovery for Deep-Koopman Operators with Kolmogorov-Arnold Networks for Dynamics}
\author{Kevin Slote, Erik Bollt, Jeremie Fish}
\date{Source: arXiv:2605.06000v1 (CC BY 4.0)}
\begin{document}
\maketitle

\noindent\emph{Source and licence.} Deep-Koopman-KANDy: Dictionary Discovery for Deep-Koopman Operators with Kolmogorov-Arnold Networks for Dynamics. Version arXiv:2605.06000v1 (CC BY 4.0), distributed by arXiv under the Creative Commons Attribution 4.0 International licence, http://creativecommons.org/licenses/by/4.0/, as recorded in the arXiv licence metadata for this exact version and shown on the abstract page https://arxiv.org/abs/2605.06000v1. Full licence notice: \texttt{licenses/arxiv-2605.06000-CC-BY-4.0.txt}.\par\smallskip

\noindent\textit{Editorial scope.} Counted unit: the article's theorem thm:intrinsic-fact (source0.tex lines 1349--1368). The article's complete original proof of it is delivered with it (source0.tex lines 1370--1397, one \texttt{proof} environment of 28 lines). No mathematics is written, repaired or re-derived: the statement and the proof are the article's own lines, in the article's own notation, and no retained line is substituted or re-flowed.  No further source-local context is needed: every cross-reference of every retained line resolves inside the two delivered spans.  Nothing was omitted from any retained span, so the package carries no omission record.  No retained line contains a citation, so the wrapper carries no bibliography and no bibliography entry.  No retained line contains a figure environment, an image-inclusion command or drawing code, so no image is delivered and the package declares no asset dependency.  Editorial formatting only, in addition: an AMS \texttt{amsart} preamble that restates the article's own declarations and macros used by the retained lines, namely the environments \texttt{theorem} on separate counters, together with the standard packages those lines need.  The retained environments print in this document as Theorem 1 in order of appearance, whereas the article prints the same items with the numbering of its own full document.  The complete native source of the article as distributed in the arXiv e-print for this exact version is bundled unchanged as \texttt{sources/199989/source0.tex}.
\begin{theorem}[Local intrinsic factorization]
\label{thm:intrinsic-fact}
The following are equivalent at each $x\in\mathcal{A}$:
\begin{enumerate}
  \item[(i)] There exist a neighborhood $U\subseteq\mathcal{A}$ of $x$ and
    $h\in C^1\bigl(g(U)\bigr)$ such that $f|_U = h\circ g|_U$.
  \item[(ii)] $\ker\!\bigl(dg_x|_{T_x\mathcal{A}}\bigr)\subseteq
    \ker\!\bigl(df_x|_{T_x\mathcal{A}}\bigr)$.
  \item[(iii)] $\nabla^{\!\mathcal{A}} f(x)$ is parallel to
    $\nabla^{\!\mathcal{A}} g(x)$.
\end{enumerate}
When these hold, the local outer derivative is given by the
\emph{intrinsic chain rule}
\begin{equation}
\label{eq:intrinsic-hprime}
  h'\bigl(g(x)\bigr) \;=\;
    \frac{\langle\nabla^{\!\mathcal{A}} f(x),\,\nabla^{\!\mathcal{A}} g(x)\rangle}
         {\|\nabla^{\!\mathcal{A}} g(x)\|^2}.
\end{equation}
\end{theorem}

\begin{proof}
\textit{(ii)\,$\Leftrightarrow$\,(iii).} As $d$-dimensional 1-forms on
$T_x\mathcal{A}$, $df_x|_{T_x\mathcal{A}}$ and $dg_x|_{T_x\mathcal{A}}$
are linearly dependent iff their kernels are nested. Since $dg_x$ has
kernel of dimension $d-1$ on $T_x\mathcal{A}$ by the assumption
$\nabla^{\!\mathcal{A}} g \neq 0$, the inclusion in (ii) is equality, and
linear dependence is equivalent (under the metric identification) to
parallelism of the gradient vectors.

\textit{(i)\,$\Rightarrow$\,(ii).} Differentiating $f|_U =
h\circ g|_U$ along any $v\in T_x\mathcal{A}$ gives
$df_x(v) = h'(g(x))\,dg_x(v)$, so $dg_x(v)=0$ implies $df_x(v)=0$.

\textit{(ii)\,$\Rightarrow$\,(i).} Since $\nabla^{\!\mathcal{A}} g(x)\neq 0$,
the implicit function theorem on $\mathcal{A}$ provides a chart
$(g,\xi_2,\dots,\xi_d)$ on a neighborhood $U$ of $x$. In these coordinates
the kernel of $dg$ is spanned by $\partial_{\xi_2},\dots,\partial_{\xi_d}$;
condition (ii) then reads $\partial f/\partial\xi_i\equiv 0$ for $i\ge 2$
on $U$. Hence $f|_U = h(g)$ with $h$ a $C^1$ function of one variable.

\textit{Formula~\eqref{eq:intrinsic-hprime}.} Take
$v=\nabla^{\!\mathcal{A}} g(x)$ in $df_x(v) = h'(g(x))\,dg_x(v)$. Since
$\nabla^{\!\mathcal{A}} g$ is tangential, $df_x(v) =
\langle\nabla f, \nabla^{\!\mathcal{A}} g\rangle =
\langle\nabla^{\!\mathcal{A}} f, \nabla^{\!\mathcal{A}} g\rangle$, and
$dg_x(v) = \|\nabla^{\!\mathcal{A}} g\|^2$. Dividing yields
\eqref{eq:intrinsic-hprime}.
\end{proof}

\end{document}
